The study of programming languages is ultimately the study of a class of \textbf{derivation systems}.\footnote{
  In some settings, these are called \textbf{formal systems} or \textbf{deductive systems}, though it's my sense that there is no real consensus as to what exactly these terms mean.
}
The programming language itself is just an \textit{implementation} of its underlying derivation systems.
As such, we need to concern ourselves at least briefly with these notions.

Let's begin with a picture.
When we program we type a bunch of symbols which get saved to a file.
Of the sequences of symbols we can type, only a handful of them will constitute sequences of \textit{tokens} in our programming language.
Of the sequences of tokens we can produce, only a handful of those will constitute \textit{well-formed} programs.
Of the well-formed programs, only a handful make sense as programs; said another way, only a handful are \textit{well-typed} programs.
Of the well-typed programs, only a hand full of \textit{those} make sense as computations; said another way, only a handful can successfully \textit{evaluated}.

\todo{A picture, nested squares}

As noted in \Cref{ch:intro}, our goal is to transform a sequence of symbols into the \textit{value} of the underlying program.
This is done by successive passes or transformations of the given data.
And at each step, there's a derivation system which formally describes which inputs are \enquote{good} and which are \enquote{garbage.}

\todo{Describe the outline of the chapter.}



We're interested in formal systems of this form.  In broad strokes, a formal system is made up of a formal language, determined by a \textbf{formal grammar} and one or more \textbf{deductive systems}.
At each level, we have a deductive system which describes a subsection of system.
If each step is \textit{decidable} or \textit{semi-decidable}, meaning we can determine which satisfy what, then we can implement it as a programming language.

\section{Inference Rules}

Deductive systems are formalizations of deductive reasoning over what are called \textbf{judgments}.
They are made up of \textbf{inference rules}, which can be though of as functions which define when we can derive a new judgment from a collection of old ones.

\begin{definition}
  A \textbf{inference rule} is...
\end{definition}

\begin{example}

\end{example}

When we define an inference rule, we're often defining a \textit{collection} of inference rules simultaneously.
This is done by the use of \textbf{meta-variables}.
A meta-variable stands for a syntactic construct whose form is determined by context.

One we have a notion of inference rule, we need the notion of a \textbf{derivation}.

\begin{definition}
A \textbf{derivation} is a tree with the following properties.
\end{definition}

\section{Derivations}
